% !TEX root = ../matan.tex

\section{Бесконечные произведения, примеры.}

\begin{Definition}{Бесконечное произведение}{}
	\textbf{Бесконечным произведением} называется предел частичных произведений
	\[
	\prod_{k=1}^{\infty} b_k := \lim_{N \to \infty} \prod_{k=1}^{N} b_k.
	\]
	Произведение называется \textbf{сходящимся}, если этот предел существует, конечен и \emph{отличен от нуля}; иначе~--- расходящимся. (Требование ненулевого предела стандартно: иначе, например, обнуление одного множителя делало бы «сходящимся» любое произведение.)
\end{Definition}

\begin{Statement}{Критерий через логарифмы}{}
	Пусть $b_k > 0$ начиная с некоторого номера $n$. Тогда
	\[
	\prod_{k=1}^{\infty} b_k \text{ сходится} \iff \sum_{k=n}^{\infty} \ln b_k \text{ сходится.}
	\]
\end{Statement}

\begin{proof}
	Для частичных произведений с номерами $N \ge n$
	\[
	\ln \prod_{k=n}^{N} b_k = \sum_{k=n}^{N} \ln b_k.
	\]
	Логарифм непрерывен и строго монотонен на $(0, +\infty)$, а экспонента~--- обратная к нему непрерывная функция. Поэтому существование конечного предела суммы $\sum \ln b_k$ равносильно существованию конечного предела произведения $\prod_{k=n}^{N} b_k$, причём предел произведения положителен (экспонента конечного числа), то есть ненулевой. Множители с номерами $k < n$ дают конечный ненулевой множитель и сходимости не меняют.
\end{proof}

\begin{Definition}{Абсолютная сходимость}{}
	Произведение $\prod_{k=1}^{\infty} b_k$ (с $b_k > 0$ начиная с некоторого номера) называется \textbf{абсолютно сходящимся}, если сходится ряд $\sum_k |\ln b_k|$.
\end{Definition}

\begin{Corollary}{Произведения вида $\prod(1 + a_k)$}{}
	Пусть $a_k > -1$ для всех $k$. Тогда произведение $\prod_{k=1}^{\infty} (1 + a_k)$ абсолютно сходится тогда и только тогда, когда сходится ряд $\sum_{k=1}^{\infty} |a_k|$.
\end{Corollary}

\begin{proof}
	Ключевая оценка: для $|t| < 1/2$
	\[
	\frac{1}{2} \le \left| \frac{\ln(1 + t)}{t} \right| \le \frac{3}{2}.
	\]
	($\Leftarrow$) Если $\sum |a_k| < \infty$, то $a_k \to 0$, поэтому при больших $k$ выполнено $|a_k| < 1/2$, и по оценке $\tfrac12 |a_k| \le |\ln(1 + a_k)| \le \tfrac32 |a_k|$. Значит, $\sum |\ln(1 + a_k)|$ сходится (мажорируется $\tfrac32 \sum |a_k|$), то есть произведение абсолютно сходится.

	($\Rightarrow$) Если произведение абсолютно сходится, то $\sum |\ln(1 + a_k)| < \infty$, откуда $\ln(1 + a_k) \to 0$ и $a_k \to 0$. При больших $k$ снова $|a_k| < 1/2$, и та же двусторонняя оценка даёт $|a_k| \le 2|\ln(1 + a_k)|$, поэтому $\sum |a_k|$ сходится. Конечное число начальных членов на сходимость не влияет.
\end{proof}

\begin{Example}{Произведение Эйлера для синуса}{}
	\[
	\sin x = x \prod_{k=1}^{\infty} \left( 1 - \frac{x^2}{\pi^2 k^2} \right), \qquad x \in \RR.
	\]
	При фиксированном $x$ члены имеют вид $1 + a_k$ с $a_k = -x^2/(\pi^2 k^2)$, и $\sum_k |a_k| = \dfrac{x^2}{\pi^2}\sum_k \dfrac{1}{k^2} < \infty$, поэтому произведение абсолютно сходится.
\end{Example}

\begin{Example}{Гамма-функция (формула Вейерштрасса)}{}
	\[
	\Gamma(s) = \frac{1}{s\, e^{\gamma s}} \prod_{k=1}^{\infty} \left( 1 + \frac{s}{k} \right)^{-1} e^{s/k}, \qquad s \neq 0, -1, -2, \ldots,
	\]
	где $\gamma = \lim_{n \to \infty}\bigl( \sum_{k=1}^{n} \tfrac{1}{k} - \ln n \bigr)$~--- постоянная Эйлера. Множители $\bigl(1 + \tfrac{s}{k}\bigr)^{-1} e^{s/k}$ ведут себя как $1 + O(1/k^2)$, что и обеспечивает сходимость произведения. Эквивалентное интегральное представление:
	\[
	\Gamma(s) = \int_0^{\infty} x^{s-1} e^{-x}\,dx \qquad (\operatorname{Re} s > 0).
	\]
\end{Example}
